\section{Experimental design}

\subsection{Tasks and controlled information regimes}

The core experiment should manipulate whether corrective evidence is available while
holding the surface task fixed. Three agentic task families provide increasing horizon:

\begin{enumerate}
  \item \textbf{Tool selection:} an agent chooses among tools after an invalid call.
  Available-tool declarations and error messages can be shown, retrievable, or hidden.
  \item \textbf{Program repair:} an agent fixes a repository after a test failure. The
  failing assertion, stack trace, or hidden-test diagnosis can be independently exposed.
  \item \textbf{Long-horizon recovery:} a coding agent receives a localized tool,
  planning, or context management failure inside a multi-step rollout and must diagnose
  it before continuing.
\end{enumerate}

Every task instance is rendered in three regimes: \emph{visible evidence}, \emph{evidence
obtainable through an allowed action}, and \emph{evidence unavailable}. This prevents a
method from receiving credit for memorizing the majority hidden answer. Counterfactual
pairs share the same visible prefix but differ in one causal variable and therefore require
different hints and actions.

\subsection{Baselines}

\begin{center}
\begin{tabular}{p{0.25\textwidth}p{0.64\textwidth}}
  \toprule
  \textbf{Method} & \textbf{What it tests} \\
  \midrule
  Direct SFT & Clone corrected actions without an explicit diagnosis. \\
  OPSD & Pull the unhinted student distribution toward a privileged
  hint-conditioned teacher \citep{shenfeld2026selfdistillation,cursor2026composer25}. \\
  Reward-only RL & Learn from task success without semantic localization. \\
  HSFT-H & Train only the explicit hint generator. \\
  HSFT-HC & Train both hint and corrected continuation. \\
  HSFT-H+OPD/RL & Generate the hint explicitly, then improve the on-policy continuation. \\
  \bottomrule
\end{tabular}
\end{center}

\subsection{Metrics and falsification tests}

\begin{enumerate}
  \item \textbf{Task success:} verified completion under visible, retrievable, and hidden
  evidence.
  \item \textbf{Hint grounding:} fraction of hint claims entailed by the visible state or
  an allowed retrieval result.
  \item \textbf{Counterfactual sensitivity:} when one evidence variable changes, does the
  hint and action change in the required direction?
  \item \textbf{Unsupported-action rate:} how often does the model confidently choose a
  correction whose required evidence is absent?
  \item \textbf{Causal mediation:} remove or shuffle the generated hint before the
  continuation. If performance does not fall, the model is not using the claimed
  intermediate.
  \item \textbf{Long-horizon stability:} success and hint accuracy as the distance from
  evidence to correction increases.
  \item \textbf{Retention and efficiency:} old-task regression, tokens generated, and
  optimizer tokens per acquired correction.
\end{enumerate}

\paragraph{Primary hypothesis.}
Under visible or retrievable evidence, hinted SFT will reduce unsupported corrections and
improve counterfactual sensitivity relative to OPSD while retaining more prior capability
than full continuation SFT. Under unavailable evidence, a properly gated hinted-SFT model
should abstain or retrieve rather than appear accurate through answer leakage.

\paragraph{Falsification criteria.}
The proposal is weakened if OPSD matches hinted SFT on counterfactual pairs, hint removal
does not reduce hinted-SFT performance, or explicit hints increase error propagation
without improving localization. It is rejected as a causal account if the model succeeds
in the hidden-evidence regime only because of dataset imbalance or contamination.
